  %%!TEX root = ./UserManual.tex
\chapter{Introduction}
\label{chap:Introduction}
  
This manual provides a brief description of the Stride software and its
features. 
Stride stands for \textbf{S}imulation of \textbf{tr}ansmission 
 of \textbf{i}nfectious \textbf{d}is\textbf{e}ases. It is an agent-based modeling 
system for close-contact infectious disease outbreaks developed by researchers at the
University of Antwerp and Hasselt University, Belgium.
The simulator uses census-based synthetic populations
that capture the demographic and geographic distributions, as well as detailed social networks.

Stride is an open source software. License information can be found in the project root directory in the file \texttt{LICENSE.txt}. The authors hope to make large-scale
agent-based epidemic models more useful to the community.

The software was developed through collaborative efforts with numerous members at various stages of its development. In alphabetical order, the project acknowledges
Philippe Beutels, Jan Broeckhove, Alexandra Cimpean, Niel Hens, Elise Kuylen, Pieter Libin, Signe Mogelmose, Astrid Sierens and Lander Willem.


More info on the project and results obtained with the software
can be found in
\begin{itemize}
\item Willem L, Stijven S, Tijskens E, Beutels P, Hens N \& Broeckhove J. (2015) Optimizing agent-based transmission models for infectious diseases, BMC Bioinformatics, 16:183 \cite{willem2015}
 
\item Kuylen E, Stijven, S, Broeckhove,J, \& Willem, L (2017) Social Contact Patterns in an Individual-based Simulator for the Transmission of Infectious Diseases, Procedia Computer Science, 108C:2438 \cite{kuylen2017}

Kuylen, E., Willem, L., Hens, N., \& Broeckhove, J. (2019). Future ramifications of age-dependent immunity levels for measles: Explorations in an individual-based model. In International conference on computational science (pp. 456–467). Springer \cite{kuylen2019future}
 
\item Kuylen, E., Willem, L., Broeckhove, J., Beutels, P., \& Hens, N. (2020). Clustering of susceptible individuals within households can drive measles outbreaks: An individual-based model exploration. Sci Rep, 10(19645). \cite{kuylen2020clustering}

\item Willem, L., Abrams, S., Libin, P., Coletti, P., Kuylen, E., Petrof, O., ... Hens, N. (2021). The impact of contact tracing and household bubbles on deconfinement strategies for COVID-19. Nature Communications, 12, 1524. \cite{willem2021impact}

\item Libin, P. J. K., Willem, L., Verstraeten, T., Torneri, A., Vanderlocht, J., \& Hens, N. (2021). Assessing the feasibility and effectiveness of household-pooled universal testing to control COVID-19 epidemics. PLoS Comp Biol, 17(3), e1008688. \cite{libin2021assessing}

\item Kuylen, E. J., Torneri, A., Willem, L., Libin, P. J. K., Abrams, S., Coletti, P., ... Hens, N. (2022). Different forms of superspreading lead to different outcomes: Heterogeneity in infectiousness and contact behavior relevant for the case of sars-cov-2. PLOS Computational Biology, 18(8), 1–20. \cite{kuylen2022different}

\end{itemize}

Please note that functionality has been improved, extended and added since those publications.


